\usepackage{tikz}
\usepackage{svg}

\usetikzlibrary{shapes.geometric}
\usetikzlibrary{arrows.meta}
\usetikzlibrary{positioning}
\usetikzlibrary{calc}
\usetikzlibrary{backgrounds}
\usetikzlibrary{fit}
\usetikzlibrary{fadings}
\usetikzlibrary{decorations}
\usepgflibrary{decorations.pathmorphing}
\usepgflibrary{decorations.markings}

\renewcommand{\familydefault}{\sfdefault}
\usepackage{sansmathfonts}

\usepackage{rotating}

\definecolor{pcbcolour}{RGB}{90, 32, 90}
\definecolor{VULEDon}{RGB}{255, 128, 0}
\definecolor{VULEDoff}{RGB}{229, 207, 184}
% \colorlet{VULEDoff}{VULEDon!0.9!white}%


% for organising the multi-panel figures
% \tikzset{panel/.style={inner sep=0pt, }};
% \tikzset{panellabel/.style={anchor=north west, inner sep=0pt, font={\bfseries\sffamily}}};
\tikzstyle{panel} = [
    inner sep=0pt,
]
\tikzstyle{panellabel} = [
    anchor=north west,
    inner sep=0pt,
    font={\bfseries\sffamily},
]
\definecolor{cExcitatory}{RGB}{255, 77, 77}
\definecolor{cInhibitory}{RGB}{0, 51, 204}
\definecolor{cBGInhibitory}{RGB}{153, 194, 255}

\tikzstyle{connection} = [
    line width=1.0pt,
    preaction={draw,white,line width=2pt, line cap=round},
    line cap=round,
    ]
\tikzstyle{connectionSmall} = [
    line width=0.5pt,
    preaction={draw,white,line width=1pt, line cap=round},
    line cap=round,
    ]

\tikzstyle{inhibitory} = [
    line width=1.0pt,
    color=cInhibitory,
    preaction={draw,white,line width=2pt, line cap=round, -{Bar[scale=0.4,sep=-1.0pt]}},
    -{Bar[scale=0.4,sep=-0.5pt]},
    ]
\tikzstyle{excitatory} = [
    line width=1.0pt,
    color=cExcitatory,
    preaction={draw,white,line width=2pt, line cap=round, -{Latex[round, scale=0.6,sep=-1.8pt]}},
    -{Latex[round, scale=0.6,sep=-0.4pt]},
    ]

\tikzstyle{neuronlabel} = [
    white,
    font=\sffamily\small,
    ]

\newcommand{\ledVU}[7]{
    % macro for sketch used in overviews
    % arg1: name of nodes
    % arg2: position of node 
    % arg3: argScale
    % arg4: how many membrane lights are on
    % arg5: is spikeLED on
    % arg6: fill bkgnd
    % arg7: draw
    % %%%%%%%%%%%%% draw stuff
    \ifthenelse{ \equal{#3}{} }{
        \def\argScale{1}
    }{
        \def\argScale{#3}
    }
    \pgfmathsetmacro\halfWidth{0.4mm * \argScale}
    \pgfmathsetmacro\segmentHalfHeight{0.4mm * \argScale}
    \pgfmathsetmacro\segmentSeparation{0.3mm * \argScale}
    \def\spikeSeparation{2.5}  % in fractions of segment separation
    \def\LEDseparation{0.2mm * \argScale}

    \pgfmathsetmacro\totalHeight{\segmentHalfHeight*2*7 + \segmentSeparation * (5 + \spikeSeparation) + 2 * \LEDseparation}
    \coordinate (#1) at (#2);
    \coordinate (#1-bottom) at ([yshift=-\totalHeight]#1);
    \coordinate (#1-bottomL) at ([xshift=-\halfWidth - \LEDseparation]#1-bottom);
    \coordinate (#1-bottomR) at ([xshift=\halfWidth + \LEDseparation]#1-bottom);
    \coordinate (#1-L) at ([xshift=-\halfWidth - \LEDseparation]#1);
    \coordinate (#1-R) at ([xshift=\halfWidth + \LEDseparation]#1);
    
    \def\VULEDon{VULEDon}
    \def\VULEDoff{VULEDoff}

    \ifthenelse{ \equal{#6}{} }{ \def\optBknd{none} }{ \def\optBknd{pcbcolour} }
    \ifthenelse{ \equal{#7}{} }{ \def\optDraw{none} }{ \def\optDraw{pcbcolour} }
    \draw [
        shape=rectangle, anchor=center,
        fill=\optBknd,
        draw=\optDraw,
        line width=0.4pt * \argScale,
        % decorate, decoration={
        %     zigzag,segment length=0.5mm * \argScale, amplitude=.04mm * \argScale
        % },
        % minimum width=2mm * \argScale, minimum height=8mm * \argScale,
    ] (#1-L) rectangle (#1-bottomR);

    % define coords
    \foreach \i in {0,...,5}
    {
        \pgfmathsetmacro\height{\segmentHalfHeight*(2*\i+1)+ \segmentSeparation*\i + \LEDseparation};
        \coordinate (#1-LED\i) at ([yshift=\height]#1-bottom);
    };
    \pgfmathsetmacro\height{\segmentHalfHeight*(2*6+1) + \segmentSeparation * (5 + \spikeSeparation) + \LEDseparation};
    \coordinate (#1-LEDspike) at ([yshift=\height]#1-bottom);

    % draw stuff
    \foreach \i in {0,...,5}
    {
        \pgfmathsetmacro\iPlusOne{\i + 1};
        \ifthenelse{\i < #4}{\def\state{VULEDon}}{\def\state{VULEDoff}}
        \draw [
            fill=\state,
            draw=none,
        ] ([xshift=-\halfWidth, yshift=-\segmentHalfHeight]#1-LED\i) rectangle ([xshift=\halfWidth, yshift=\segmentHalfHeight]#1-LED\i) ;
        \ifthenelse{\i < #4}{
            \foreach \a in {0, 45, 90, 135, ..., 360}
            {
            \draw[
                rotate=\a,
                shorten <=1pt,
                draw=VULEDon,
            ] (#1-LED\i) -- ++(0.08*\argScale, 0);
            }
        }{}
    };
    \ifthenelse{ \equal{#5}{} }{\def\state{VULEDoff}}{\def\state{VULEDon}}
    \draw [
        fill=\state,
        draw=none,
    ] ([xshift=-\halfWidth, yshift=-\segmentHalfHeight]#1-LEDspike) rectangle ([xshift=\halfWidth, yshift=\segmentHalfHeight]#1-LEDspike) ;
    \ifthenelse{ \equal{#5}{} }{}{
            \foreach \a in {0, 45, 90, 135, ..., 360}
            {
            \draw[
                rotate=\a,
                shorten <=1pt,
                draw=VULEDon,
            ] (#1-LEDspike) -- ++(0.08*\argScale, 0);
            }
    }
}

\newcommand{\pcbNeuron}[7]{
    % macro for sketch used in overviews
    % arg1: name of nodes
    % arg2: point of node 
    % arg3: argScale
    % arg4: rotation
    % arg5: if nonempty, use only border shape of PCB
    % arg6: if nonempty, draw LEDs (arg6 are on, if >6 spiking)
    % arg7: if nonempty, draw tunabels
    % %%%%%%%%%%%%% draw stuff

    \ifthenelse{ \equal{#3}{} }{
        \def\argScale{1}
    }{
        \def\argScale{#3}
    }
    \ifthenelse{ \equal{#4}{} }{
        \def\rot{0}
    }{
        \def\rot{#4}
    }

    \coordinate (#1) at (#2);

    % label coord
    \coordinate (#1-l) at ([yshift=0.4mm * \argScale, xshift=0.6mm * \argScale] #2);

    \node[anchor=center,
          xshift=0.25cm * \argScale,yshift=0.1cm * \argScale, 
          rotate around={\rot:(#1)},
          ] (whitehead) at (#1)
    {
        % \scaledWidth
        \scalebox{\argScale}{
            \ifthenelse{ \equal{#5}{} }{
                \includesvg[width=2cm]{neuron_filled.svg}
            }{
                \includesvg[width=2cm]{neuron_border.svg}
            }
        }
    };

    % %%%%%%%%%%%%%%%%%%% dendrites
    \pgfmathsetmacro{\lengthDaX}{-0.09cm * \argScale}
    \pgfmathsetmacro{\lengthDaY}{0.64cm * \argScale}
    \coordinate[rotate around={\rot:(#1)}] (#1-Da) at ([xshift=\lengthDaX, yshift=\lengthDaY] #1);

    \pgfmathsetmacro{\lengthDbX}{-0.685cm * \argScale}
    \pgfmathsetmacro{\lengthDbY}{0.138cm * \argScale}
    \coordinate[rotate around={\rot:(#1)}]  (#1-Db) at ([xshift=\lengthDbX, yshift=\lengthDbY] #1);

    \pgfmathsetmacro{\lengthDcX}{-0.47cm * \argScale}
    \pgfmathsetmacro{\lengthDcY}{-0.42cm * \argScale}
    \coordinate[rotate around={\rot:(#1)}]  (#1-Dc) at ([xshift=\lengthDcX, yshift=\lengthDcY] #1);

    % %%%%%%%%%%%%%%%%%%% axons
    \pgfmathsetmacro{\lengthAaX}{1.150cm * \argScale}
    \pgfmathsetmacro{\lengthAaY}{-0.360cm * \argScale}
    \coordinate[rotate around={\rot:(#1)}]  (#1-Aa) at ([xshift=\lengthAaX, yshift=\lengthAaY] #1);

    \pgfmathsetmacro{\lengthAbX}{1.135cm * \argScale}
    \pgfmathsetmacro{\lengthAbY}{-0.390cm * \argScale}
    \coordinate[rotate around={\rot:(#1)}]  (#1-Ab) at ([xshift=\lengthAbX, yshift=\lengthAbY] #1);

    \pgfmathsetmacro{\lengthAcX}{1.120cm * \argScale}
    \pgfmathsetmacro{\lengthAcY}{-0.420cm * \argScale}
    \coordinate[rotate around={\rot:(#1)}]  (#1-Ac) at ([xshift=\lengthAcX, yshift=\lengthAcY] #1);

    % %%%%%%%%%%%%%%%%%%% LEDs
    \pgfmathsetmacro{\lengthLEDX}{0.770cm * \argScale}
    \pgfmathsetmacro{\lengthLEDY}{-0.205cm * \argScale}
    \pgfmathsetmacro{\newscaling}{\argScale * 0.8}  % isn't tex beautiful..
    \coordinate[rotate around={\rot:(#1)}] (#1-LEDs) at ([xshift=\lengthLEDX, yshift=\lengthLEDY] #1);
    \ifthenelse{ \equal{#6}{} }{}{
        \ifthenelse{ #6 > 6}{
            \node[rotate=244+\rot] (tmp) at (#1-LEDs) {
                \begin{tikzpicture}[
                ]
                \ledVU{luiLED}{0, 0}{\newscaling}{0}{yes}{}{} ;
                \end{tikzpicture}
            };
        }{
            \node[rotate=244+\rot] (tmp) at (#1-LEDs) {
                \begin{tikzpicture}[
                ]
                \ledVU{luiLED}{0, 0}{\newscaling}{#6}{}{}{} ;
                \end{tikzpicture}
            };
        };
    }

    % %%%%%%%%%%%%%%%%%%% tunables
    \pgfmathsetmacro{\lengthTunPotiRad}{1.29 * \argScale}
    \pgfmathsetmacro{\lengthTunPotiScrew}{0.59 * \argScale}

    \pgfmathsetmacro{\lengthTunVleakX}{-0.059cm * \argScale}
    \pgfmathsetmacro{\lengthTunVleakY}{0.135cm * \argScale}
    \pgfmathsetmacro{\lengthTunTaumX}{-0.059cm * \argScale}
    \pgfmathsetmacro{\lengthTunTaumY}{0.029cm * \argScale}
    \pgfmathsetmacro{\lengthTunTausX}{-0.059cm * \argScale}
    \pgfmathsetmacro{\lengthTunTausY}{-0.076cm * \argScale}
    \coordinate[rotate around={\rot:(#1)}] (#1-tun-v_leak) at ([xshift=\lengthTunVleakX, yshift=\lengthTunVleakY] #1);
    \coordinate[rotate around={\rot:(#1)}] (#1-tun-taum) at ([xshift=\lengthTunTaumX, yshift=\lengthTunTaumY] #1);
    \coordinate[rotate around={\rot:(#1)}] (#1-tun-taus) at ([xshift=\lengthTunTausX, yshift=\lengthTunTausY] #1);

    % weight potis
    \pgfmathsetmacro{\lengthTunWaX}{0.0cm * \argScale}
    \pgfmathsetmacro{\lengthTunWaY}{0.264cm * \argScale}
    \pgfmathsetmacro{\lengthTunWbX}{-0.329cm * \argScale}
    \pgfmathsetmacro{\lengthTunWbY}{0.076cm * \argScale}
    \pgfmathsetmacro{\lengthTunWcX}{-0.265cm * \argScale}
    \pgfmathsetmacro{\lengthTunWcY}{-0.129cm * \argScale}
    \coordinate[rotate around={\rot:(#1)}] (#1-tun-Wa) at ([xshift=\lengthTunWaX, yshift=\lengthTunWaY] #1);
    \coordinate[rotate around={\rot:(#1)}] (#1-tun-Wb) at ([xshift=\lengthTunWbX, yshift=\lengthTunWbY] #1);
    \coordinate[rotate around={\rot:(#1)}] (#1-tun-Wc) at ([xshift=\lengthTunWcX, yshift=\lengthTunWcY] #1);

    % drawing
    \ifthenelse{ \equal{#7}{} }{}{
        \foreach \iOtherIterator in {v_leak, taum, taus, Wa, Wb, Wc}
        {
            \fill[black!30!white] (#1-tun-\iOtherIterator) circle (\lengthTunPotiRad pt) coordinate (tmp);
            \draw[black!60!white, rotate around={\rot:(tmp)}] (tmp) ++ (-\lengthTunPotiScrew pt, 0) -- ++(2 * \lengthTunPotiScrew pt, 0);
            \draw[black!60!white, rotate around={\rot:(tmp)}] (tmp) ++ (0, -\lengthTunPotiScrew pt) -- ++(0, 2 * \lengthTunPotiScrew pt);
        };
    }

    % weight switches
    \pgfmathsetmacro{\lengthTunWSaX}{-0.082cm * \argScale}
    \pgfmathsetmacro{\lengthTunWSaY}{0.394cm * \argScale}
    \pgfmathsetmacro{\lengthTunWSbX}{-0.412cm * \argScale}
    \pgfmathsetmacro{\lengthTunWSbY}{0.153cm * \argScale}
    \pgfmathsetmacro{\lengthTunWScX}{-0.329cm * \argScale}
    \pgfmathsetmacro{\lengthTunWScY}{-0.276cm * \argScale}
    \coordinate[rotate around={\rot:(#1)}] (#1-tun-WSa) at ([xshift=\lengthTunWSaX, yshift=\lengthTunWSaY] #1);
    \coordinate[rotate around={\rot:(#1)}] (#1-tun-WSb) at ([xshift=\lengthTunWSbX, yshift=\lengthTunWSbY] #1);
    \coordinate[rotate around={\rot:(#1)}] (#1-tun-WSc) at ([xshift=\lengthTunWScX, yshift=\lengthTunWScY] #1);

    % drawing
    \ifthenelse{ \equal{#7}{} }{}{
        \foreach \iOtherIterator / \angleUnrotated in {WSa/81, WSb/4, WSc/60}
        {
            \pgfmathsetmacro{\angleRotated}{\angleUnrotated + \rot}
            \fill[black!90!white, rotate around={\angleRotated:(#1-tun-\iOtherIterator)}] (#1-tun-\iOtherIterator) ++(-0.024*\argScale, 0.0) rectangle ++(0.047*\argScale, 0.118*\argScale);
            \fill[black!30!white, rotate around={\angleRotated:(#1-tun-\iOtherIterator)}] (#1-tun-\iOtherIterator) ++(-0.079*\argScale, 0.0) rectangle ++(0.159*\argScale, 0.059*\argScale);
        };
    }
}
